\chapter{Results and Probabilistic Forecasts}
\label{chapter:results}

This chapter presents the quantitative outputs generated by aggregating the individual expert probability distributions into a single synthetic expert distribution using the performance-maximized Decision Maker ($\text{PWDM}_{\alpha^*}$). For each variable of interest, the subjective uncertainties are expressed via the elicited 5\%, 50\%, and 95\% quantiles to establish a robust, non-deterministic 2027 residential market forecast.

\section{Target Question 1: Total Newly Completed Dwellings in 2027}
\label{sec:target_q1}

The optimal Decision Maker returns the following aggregated subjective quantiles for the total number of newly constructed residential units realized across the Netherlands during the calendar year 2027: ($q_{0.05} = [\text{X}]$, $q_{0.50} = [\text{Y}]$, $q_{0.95} = [\text{Z}]$).\\\newline
Therefore, the mathematically optimized median estimate for newly completed national dwellings in 2027 is $[\text{Y}]$ units. Furthermore, the expert panel establishes a 90\% subjective probability that the true realized construction output will fall within the continuous interval $[[\text{X}], [\text{Z}]]$. This distribution formally quantifies the structural supply-side headwinds and labor constraints previously identified in the market analysis.

\section{Target Question 2: Total Residential Building Permits Issued in 2027}
\label{sec:target_q2}

The optimal Decision Maker returns the following aggregated subjective quantiles for the absolute number of municipal residential building permits formally granted during the calendar year 2027: ($q_{0.05} = [\text{X}]$, $q_{0.50} = [\text{Y}]$, $q_{0.95} = [\text{Z}]$).\\\newline
Consequently, the performance-weighted median expectation for pipeline permit allocations is fixed at $[\text{Y}]$. The panel expresses a 90\% subjective confidence interval bounding this indicator between $[\text{X}]$ and $[\text{Z}]$, providing a probabilistic baseline to evaluate whether upcoming project pipelines will catch up to long-term national policy milestones.

\section{Target Question 3: Average Real Estate Transaction Price in Amsterdam in 2027}
\label{sec:target_q3}

The optimal Decision Maker returns the following aggregated subjective quantiles for the average transactional selling price (in Euros) of an existing owner-occupied residential home within the municipal boundary of Amsterdam over the entire calendar year 2027: ($q_{0.05} = \text{\EUR}[\text{X}]$, $q_{0.50} = \text{\EUR}[\text{Y}]$, $q_{0.95} = \text{\EUR}[\text{Z}]$).\\\newline
The performance-weighted median prediction indicates an average transaction price of $\text{\EUR}[\text{Y}]$, with the joint expert distribution reflecting a 90\% probability that the true economic realization will be contained within the interval $[\text{\EUR}[\text{X}], \text{\EUR}[\text{Z}]]$. This spread explicitly highlights the panel's combined uncertainty regarding urban affordability trajectories and localized demand pressures.

\section{Target Question 4: Free-Sector Rent Escalation Coefficients}
\label{sec:target_q4}

The optimal Decision Maker returns the following aggregated subjective quantiles for the annualized percentage increase in average free-sector rental rates (\textit{vrijesectorwoningen}) across the Netherlands as of July 1, 2027, relative to indices from July 1, 2026: ($q_{0.05} = [\text{X}]\%$, $q_{0.50} = [\text{Y}]\%$, $q_{0.95} = [\text{Z}]\%$).\\\newline
The performance-weighted median expectation for free-sector rental appreciation is bounded at $[\text{Y}]\%$. The resulting 90\% subjective confidence interval spans from $[\text{X}]\%$ to $[\text{Z}]\%$, capturing the latent volatility introduced by shifting legislative framework controls and unresolved housing tightness.

